\documentclass{article}
\usepackage[UTF8]{ctex}
\usepackage{amsmath,amssymb,amsfonts}
\usepackage{enumitem}
\usepackage{geometry}
\geometry{a4paper, margin=1in}
% === 错题本排版优化（不参与题目标识） ===
\setlength{\parindent}{0pt}
\setlength{\parskip}{0.35em plus 0.10em minus 0.08em}
\setlength{\emergencystretch}{3em}
\setlist{itemsep=0.25em, topsep=0.40em, parsep=0pt}
\allowdisplaybreaks[2]
% === 错题本排版优化结束 ===
\begin{document}
\section*{西南财经大学2026年高等代数真题}
\section*{166．西南财经大学2026年高等代数试卷}

\subsection{题目1}
1．求行列式

\[
\left|\begin{array}{ccccc}
a & b & b & b & b \\
b & a & b & b & b \\
b & b & a & b & b \\
b & b & b & a & b \\
b & b & b & b & a
\end{array}\right|
\]

\subsection{题目2}
2．已知 \(A^{*}=\left(\begin{array}{cccc}1 & 1 & -1 & 1 \\ -1 & 1 & 1 & 1 \\ 1 & -1 & 1 & 1 \\ 0 & 0 & 0 & -2\end{array}\right)\) ，且 \(A^{*} B A=2 B A+E\) ．

\begin{itemize}
  \item[（1）] 求 \(|A|\) 。
  \item[（2）] 求矩阵 \(B\) 。
\end{itemize}

\subsection{题目3}
3．已知

\[
\beta_{1}=\alpha_{2}+\alpha_{3}+\cdots+\alpha_{s}, \beta_{2}=\alpha_{1}+\alpha_{3}+\cdots+\alpha_{s}, \cdots, \beta_{s}=\alpha_{1}+\alpha_{2}+\cdots+\alpha_{s-1} .
\]

证明向量组 \(\alpha_{1}, \alpha_{2}, \cdots, \alpha_{a}\) 与 \(\beta_{1}, \beta_{2}, \cdots, \beta_{a}\) 的秩相同。

\subsection{题目4}
4．已知三元实二次型 \(f(X)=X^{\mathrm{T}} A X\) 在正交线性替换 \(X=Q Y\) 下化为标准形 \(f=y_{2}^{2}+y_{3}^{2}\) ，其中 \(Q\)的第一列为 \(\left(\frac{\sqrt{2}}{2}, 0, \frac{\sqrt{2}}{2}\right)^{\mathrm{T}}\) ，求矩阵 \(A\) 。

\subsection{题目5}
5．求以 \(\sqrt{2}+\sqrt{3}\) 为根的有理数域上的不可约多项式．

\subsection{题目6}
6．已知

\[
\begin{aligned}
& \alpha_{1}=(1,2,1,0)^{\mathrm{T}}, \alpha_{2}=(-1,1,1,1)^{\mathrm{T}} \\
& \beta_{1}=(2,-1,0,1)^{\mathrm{T}}, \beta_{2}=(1,-1,3,7)^{\mathrm{T}}
\end{aligned}
\]

且 \(V_{1}=L\left(\alpha_{1}, \alpha_{2}\right), V_{2}=L\left(\beta_{1}, \beta_{2}\right)\) ，求 \(V_{1} \cap V_{2}\) 的基与维数。

\subsection{题目7}
7．设 \(\alpha_{1}, \alpha_{2}, \alpha_{3}, \alpha_{4}\) 是数域 \(P\) 上4维线性空间 \(V\) 的一组基，线性变换 \(\mathcal{T}\) 在这组基下的矩阵为

\[
A=\left(\begin{array}{cccc}
-1 & 3 & 4 & 2 \\
-2 & 5 & 1 & 3 \\
-6 & 17 & 9 & 11 \\
-7 & 18 & -17 & 11
\end{array}\right)
\]

\begin{itemize}
  \item[（1）] 求 \(\mathcal{T}\) 的核及其一组基。
  \item[（2）] 求 \(\mathcal{T}\) 的值域及其一组基。
\end{itemize}

\subsection{题目8}
8．设 \(\eta\) 是 \(n\) 维欧氏空间 \(V\) 的一个单位向量，定义变换 \(\sigma\) 为 \(\sigma(\alpha)=\alpha-2(\alpha, \eta) \eta\) 。

\begin{itemize}
  \item[（1）] 证明：\(\sigma\) 是正交变换．
  \item[（2）] 设 \(\sigma\) 在 \(V\) 的任意一组标准正交基下的矩阵为 \(A\) ，证明：\(|A|=-1\) 。
\end{itemize}
\end{document}
